% Compilar desde la raíz del repositorio.
% Versión compacta: 33 diapositivas (22 originales + 50%).
\input{config.tex}
\begin{document}
\begin{frame}{Funciones de varias variables}
\centering\LARGE Tema 1\par\bigskip
\large Matemáticas 2\par Grado en Ingeniería Informática
\end{frame}
\begin{frame}{Una y dos variables}
Una función asigna una salida a cada entrada admitida.
\[f(x)=x^2,\qquad f(2)=4.\]
Con dos variables, la entrada es un par ordenado:
\[g(x,y)=x+y,\qquad g(2,3)=5.\]
En ambos casos, cada entrada tiene una única salida.
\end{frame}
\begin{frame}{Temperatura terrestre}
A una hora determinada, la temperatura en la superficie de la Tierra depende de la longitud y la latitud.
\[T=f(x,y).\]
$x$ representa la longitud e $y$ la latitud.
\medskip
La función asigna una temperatura a cada posición admitida en el modelo.
\end{frame}
\begin{frame}{Dos entradas}
\[\underbrace{(x,y)}_{\text{entrada}}\ \longmapsto\ \underbrace{f(x,y)}_{\text{salida}}\]
Un par ordenado contiene dos datos y constituye una entrada.
\[f(x,y)=x+y,\qquad f(2,3)=2+3=5.\]
La salida sigue siendo un número real.
\end{frame}
\begin{frame}{Evaluación}
Para evaluar una función, sustituimos cada variable por su valor.
\[f(x,y)=xy,\qquad f(2,3)=6,\qquad f(0,4)=0.\]
El orden de las entradas importa:
\[g(x,y)=x+2y,\qquad g(1,3)=7,\qquad g(3,1)=5.\]
\textbf{Práctica:} calcula $g(2,1)$ y $g(1,2)$.
\end{frame}
\begin{frame}{Definición}
Una función real de dos variables reales asigna a cada par $(x,y)$ de un conjunto $D\subseteq\mathbb R^2$ un único número real.
\[f:D\longrightarrow\mathbb R,\qquad (x,y)\longmapsto f(x,y).\]
$\mathbb R^2$ es el conjunto de todos los pares de números reales.
\medskip
El conjunto $D$ se llama \textbf{dominio}.
\end{frame}
\begin{frame}{Variables}
\[z=f(x,y)\]
\begin{itemize}
\item $x$ e $y$: variables independientes, las entradas del modelo.
\item $z$: variable dependiente, la salida calculada.
\end{itemize}
En $T=f(x,y)$, la temperatura depende de la posición.
\medskip
Las entradas deben respetar las condiciones del dominio.
\end{frame}
\begin{frame}{Dominio e imagen}
\[f:\mathbb R^2\to\mathbb R,\qquad f(x,y)=x^2+y^2.\]
\textbf{Dominio}: entradas admitidas, $D=\mathbb R^2$.
\medskip
\textbf{Imagen}: salidas obtenidas, $\operatorname{Im}(f)=[0,\infty)$.
\medskip
La función nunca devuelve un valor negativo. Cualquier $k\ge0$ se obtiene con $(\sqrt k,0)$.
\medskip
El codominio declarado es $\mathbb R$ y contiene a la imagen.
\end{frame}
\begin{frame}{Restricciones conocidas}
Trabajamos con valores reales.
\begin{center}\begin{tabular}{ll}
Expresión & Condición\\\hline
$1/x$ & $x\ne0$\\[3mm]
$\sqrt{x}$ & $x\ge0$\\[3mm]
$\ln x$ & $x>0$\\[3mm]
$\arcsin x$ & $-1\le x\le1$
\end{tabular}\end{center}
En varias variables aplicamos las mismas condiciones a la expresión completa del denominador o del argumento.
\end{frame}
\begin{frame}{Una restricción}
\[f(x,y)=\frac1{x-y}\]
El denominador debe ser distinto de cero.
\[x-y\ne0\quad\Longleftrightarrow\quad y\ne x.\]
\[D=\{(x,y)\in\mathbb R^2:y\ne x\}.\]
Admitimos todo el plano excepto la recta $y=x$.
\end{frame}
\begin{frame}{Dominio parabólico}
\[f(x,y)=\frac{3x-5y}{y-x^2}\]
El denominador exige $y\ne x^2$.
\[D=\{(x,y)\in\mathbb R^2:y\ne x^2\}.\]
\figura{s10_4.png}{.65}{.36}
Excluimos la parábola $y=x^2$. Un numerador nulo está permitido si el denominador no es cero.
\end{frame}
\begin{frame}{Dominio de una raíz}
\[f(x,y)=\sqrt{9-x^2-y^2}\]
\[9-x^2-y^2\ge0\quad\Longleftrightarrow\quad x^2+y^2\le9.\]
El dominio es el disco de centro $(0,0)$ y radio $3$, incluida la frontera.
\medskip
$f(0,0)=3$ y $f(3,0)=0$. El punto $(4,0)$ queda fuera, pues el radicando sería negativo.
\end{frame}
\begin{frame}{Tres variables}
El volumen de una caja depende de largo, ancho y alto.
\[V(x,y,z)=xyz,\qquad V(2,3,4)=24.\]
Si las longitudes están en metros, el volumen está en $\mathrm{m}^3$.
\figura{s3_2.png}{.40}{.35}
\end{frame}
\begin{frame}{$n$ variables}
Una $n$-tupla es una lista ordenada de $n$ números:
\[(x_1,x_2,\ldots,x_n)\in\mathbb R^n.\]
Una función real de $n$ variables reales asigna a cada $n$-tupla de su dominio un único número real.
\[f:D\subseteq\mathbb R^n\longrightarrow\mathbb R,\qquad u=f(x_1,\ldots,x_n).\]
\end{frame}
\begin{frame}{Dominio logarítmico}
\[f(x,y,z)=\ln(z-y)+xy\sin z\]
El logaritmo exige $z-y>0$. El término $xy\sin z$ está definido para toda terna real.
\[D=\{(x,y,z)\in\mathbb R^3:z>y\}.\]
Es el semiespacio por encima del plano $z=y$, sin incluir el plano.
\end{frame}
\begin{frame}{Suma, producto y cociente}
Sean $f(x,y)=x+y$ y $g(x,y)=xy$.
\[(f+g)(x,y)=x+y+xy,\]
\[(fg)(x,y)=(x+y)xy,\]
\[\left(\frac fg\right)(x,y)=\frac{x+y}{xy}.\]
La suma y el producto requieren entradas de ambos dominios.
\medskip
El cociente exige además $g(x,y)\ne0$: aquí, $x\ne0$ e $y\ne0$.
\end{frame}
\begin{frame}{Composición}
La salida de una función puede ser la entrada de otra.
\[g(x,y)=x+y,\qquad h(t)=t^2.\]
\[(h\circ g)(x,y)=h(g(x,y))=(x+y)^2.\]
Para componer, $g(x,y)$ debe pertenecer al dominio de $h$.
\medskip
Si $h(t)=\sqrt t$, la composición exige $x+y\ge0$.
\end{frame}
\begin{frame}{Dimensiones de la gráfica}
\begin{center}\begin{tabular}{lll}
Función & Punto de la gráfica & Espacio\\\hline
$y=f(x)$ & $(x,f(x))$ & $\mathbb R^2$\\[3mm]
$z=g(x,y)$ & $(x,y,g(x,y))$ & $\mathbb R^3$
\end{tabular}\end{center}
Para $g(x,y)=x+y$, la entrada $(1,2)$ produce el punto $(1,2,3)$. La gráfica es un plano.
\medskip
En tres variables, $(x,y,z,f(x,y,z))\in\mathbb R^4$. Podemos estudiar cortes o conjuntos de nivel.
\end{frame}
\begin{frame}{Gráfica de dos variables}
\begin{columns}[c]
\begin{column}{.52\textwidth}
$(x,y)$ localiza la entrada. $f(x,y)$ determina la altura.
\medskip
La gráfica reúne los puntos
\[(x,y,f(x,y)),\quad (x,y)\in D.\]
Cada entrada admitida tiene una única altura.
\end{column}
\begin{column}{.45\textwidth}\figura{s8_2.png}{1}{.63}\end{column}
\end{columns}
\end{frame}
\begin{frame}{Una superficie}
\[z=f(x,y)=x+y^2,\qquad D=\mathbb R^2.\]
\figura{s9_2.png}{.85}{.52}
Con $x=0$ aparece $z=y^2$. Con $y=0$ aparece $z=x$.
\end{frame}
\begin{frame}{Una semiesfera}
\[z=\sqrt{9-x^2-y^2}.\]
\figura{s11_5.png}{.72}{.46}
Al elevar al cuadrado: $x^2+y^2+z^2=9$, con $z\ge0$.

La gráfica es la semiesfera superior. Su dominio es un disco en el plano $xy$.
\end{frame}
\begin{frame}{Altura en un mapa}
\begin{columns}[c]
\begin{column}{.46\textwidth}
La altura de un terreno depende de la posición: $H=f(x,y)$.
\medskip
Una línea reúne puntos con igual altura. Por ejemplo:
\[H(x,y)=600.\]
Fijamos la salida y buscamos las posiciones que la producen.
\end{column}
\begin{column}{.51\textwidth}\figura{s21_0.jpg}{1}{.67}\end{column}
\end{columns}
\end{frame}
\begin{frame}{Curvas de nivel}
Para $f:D\subseteq\mathbb R^2\to\mathbb R$, el conjunto de nivel $k$ es
\[L_k=\{(x,y)\in D:f(x,y)=k\}.\]
Cuando forma una curva, hablamos de \textbf{curva de nivel}.
\medskip
Si $k\notin\operatorname{Im}(f)$, el conjunto es vacío. Algunos niveles pueden ser un solo punto.
\end{frame}
\begin{frame}{Niveles rectos}
\[f(x,y)=x+y\]
El nivel $k$ cumple $x+y=k$, es decir, $y=k-x$.
\begin{center}\begin{tabular}{ccl}
Nivel & Ecuación & Puntos del nivel\\\hline
$0$ & $y=-x$ & $(0,0)$, $(1,-1)$\\[2mm]
$2$ & $y=2-x$ & $(0,2)$, $(1,1)$\\[2mm]
$4$ & $y=4-x$ & $(0,4)$, $(2,2)$
\end{tabular}\end{center}
Las curvas de nivel son rectas paralelas.
\end{frame}
\begin{frame}{Niveles circulares}
\[f(x,y)=x^2+y^2\]
El nivel $k$ cumple $x^2+y^2=k$.
\begin{itemize}
\item $k>0$: circunferencia de radio $\sqrt k$.
\item $k=0$: únicamente el punto $(0,0)$.
\item $k<0$: conjunto vacío.
\end{itemize}
La imagen $[0,\infty)$ indica qué niveles existen.
\end{frame}
\begin{frame}{Superficie y niveles}
\begin{columns}[c]
\begin{column}{.48\textwidth}\figura{s14_3.png}{1}{.58}\end{column}
\begin{column}{.48\textwidth}\figura{s18_0.jpg}{1}{.58}\end{column}
\end{columns}
Un corte horizontal $z=k$ reúne puntos de igual altura. Su proyección al plano $xy$ muestra el nivel $f(x,y)=k$.
\end{frame}
\begin{frame}{Niveles de una raíz}
\[f(x,y)=\sqrt{100-x^2-y^2}\]
El dominio es $x^2+y^2\le100$ y la imagen es $[0,10]$.
\medskip
Buscamos las posiciones cuya salida sea $c$:
\[\sqrt{100-x^2-y^2}=c.\]
Primero debe cumplirse $0\le c\le10$.
\end{frame}
\begin{frame}{Ecuación del nivel}
Para $0\le c\le10$,
\[\sqrt{100-x^2-y^2}=c
\quad\Longleftrightarrow\quad x^2+y^2=100-c^2.\]
Con $0\le c<10$, obtenemos circunferencias de radio $\sqrt{100-c^2}$.
\medskip
Con $c=10$, el nivel es el punto $(0,0)$.

Con $c<0$ o $c>10$, es vacío.
\medskip
Exigir $c\ge0$ evita soluciones falsas al elevar al cuadrado.
\end{frame}
\begin{frame}{Radios y alturas}
\begin{columns}[c]
\begin{column}{.50\textwidth}
\[x^2+y^2=100-c^2\]
\begin{center}\begin{tabular}{cc}
Altura $c$ & Radio\\\hline
$0$ & $10$\\[2mm]
$6$ & $8$\\[2mm]
$8$ & $6$\\[2mm]
$10$ & $0$
\end{tabular}\end{center}
Al subir, disminuye el radio.
\end{column}
\begin{column}{.47\textwidth}\figura{s17_3.png}{1}{.57}\end{column}
\end{columns}
\end{frame}
\begin{frame}{Lectura de niveles}
\begin{columns}[c]
\begin{column}{.48\textwidth}\figura{s19_1.jpg}{1}{.39}\end{column}
\begin{column}{.48\textwidth}\figura{s20_1.jpg}{1}{.39}\end{column}
\end{columns}
Un mismo nivel puede tener varias componentes separadas.
\medskip
En un mapa térmico de un procesador, $T(x,y)=70$ reúne posiciones a $70\,{}^\circ\mathrm C$. En meteorología, las isobaras reúnen puntos de igual presión.
\end{frame}
\begin{frame}{Superficies de nivel}
En tres variables, fijar la salida suele dar una superficie:
\[L_k=\{(x,y,z)\in D:f(x,y,z)=k\}.\]
Si $f=x^2+y^2+z^2$, el nivel $k=4$ es una esfera de radio $2$.
\medskip
Si $f=x^2+y^2$, para $k>0$ obtenemos cilindros de radio $\sqrt k$: $z$ queda libre.

Para $k=0$, obtenemos el eje $z$. Para $k<0$, el conjunto vacío.
\end{frame}
\begin{frame}{Ejercicio de dominio}
Encuentra el dominio natural de
\[f(x,y,z)=\arcsin(xy)e^{3z}+e^{(y+z)\ln(x+z)}.\]
\begin{enumerate}
\item Recuerda: $\arcsin t$ exige $-1\le t\le1$.
\item Recuerda: $\ln t$ exige $t>0$.
\item Reúne todas las condiciones en un conjunto de ternas.
\end{enumerate}
Las exponenciales admiten cualquier exponente real definido.
\end{frame}
\begin{frame}{Ejercicios de nivel}
Describe los conjuntos de nivel:
\begin{enumerate}
\item $f(x,y)=x^2+y^2$,\quad $k=0,1,2,3$.
\item $g(x,y,z)=x^2+y^2$,\quad $k=0,1,2$.
\item $h(x,y)=x^2-y^2$,\quad $k=-2,-1,0,1,2$.
\end{enumerate}
En 1 y 3 estudiamos niveles en el plano. En 2, niveles en el espacio.
\medskip
Para $k=0$ en 3, usa $x^2-y^2=(x-y)(x+y)$.
\end{frame}
\end{document}
